\documentclass[11pt]{article}

\usepackage[margin=1in]{geometry}
\usepackage{booktabs}
\usepackage{array}
\usepackage{hyperref}
\usepackage{longtable}
\usepackage{enumitem}

\title{Comparing Data Availability in Quantum and Classical Computing Research:\\
A Single Two-Axis, Six-Case Classification Pipeline}
\author{}
\date{\today}

\begin{document}
\maketitle

\begin{abstract}
This report documents a single pipeline for comparing data-availability practices
between quantum computing (superconducting-qubit) research and classical
(``regular'') computing research: whether a paper relies on data external to itself,
on data it generates itself, or both, and how completely and verifiably each was
shared. The pipeline is defined once and configured with a set of subject-matter
pillars and keyword definitions for each side of the comparison; it builds a
per-pillar candidate corpus from the arXiv API for each side, screens every candidate
via a full-text scan for data/code-availability signal, and classifies each paper by a
two-axis, six-case framework. A separate, independently-judged classification
determines each paper's true subject pillar, correcting for the fact that the pillar
under which a paper was originally searched need not reflect what the paper's content
is actually about. This report also documents two corpus-construction defects
discovered and corrected during this study --- a silent 100-result-per-keyword
truncation in the arXiv query step, and a systematic undercount in the
classical-computing keyword list caused by generic descriptive phrases (e.g.\
``machine learning algorithm'') that do not match how the literature actually names
its own methods --- together with the fixes adopted for each. At the time of writing,
3{,}472 of 3{,}478 quantum-side papers have completed the two-axis classification; the
classical-side corpus (6{,}308 papers, pulled after the truncation fix but before the
keyword-list rework) has not yet been classified. A worked pilot sample of fourteen
papers per side, classified by hand prior to the scaled pipeline described here, is
reported in Section~\ref{sec:results-detail} as an illustration of the framework; full
results for both sides of the comparison, and the corpus-construction diagnostic
tables referenced above, will be added to this report once classification of the
corrected, full-breadth corpora is complete.
\end{abstract}

\section{Introduction}

Reproducibility audits of empirical research papers typically ask a single question ---
``is the data available?'' --- that conflates several distinct concerns: whether any
data was used at all, whether that data came from outside the paper or was produced by
the paper's own procedure, and whether what was used was actually handed to the reader
in a verifiable form. A paper can use data drawn from outside itself, data its own code
or experiment produces, both, or neither, and each of these can in turn be handed over
completely, named but not shared, shared incompletely, claimed as shared but
unreachable, or shared only conditionally on request. Collapsing all of this into one
yes/no judgment obscures which of these situations a given paper is actually in. This
report describes a single pipeline that separates those concerns into a two-axis
classification scheme --- whether external data is used, and whether self-generated
data is used, each independently scored across six possible states of availability.
The pipeline is defined once, independent of subject matter, and configured with a set
of subject-matter pillars and keyword definitions for each side of an ongoing
comparison of data-availability practices between quantum computing
(superconducting-qubit) research and classical computing research, giving one
independently-built corpus and classification output per side of that comparison. The
pipeline that generates and screens candidate papers is described in
Section~\ref{sec:candidate-generation}; the classification framework itself, and its
output on a pilot sample drawn from each side's review queue, occupy the remainder of
this report.

\section{Methodology}

\subsection{Candidate paper generation}
\label{sec:candidate-generation}

The papers reviewed in this pass were drawn from an existing, larger
candidate-generation pipeline: a single, reusable procedure, run once per side of a
broader thesis comparing data-availability practices between quantum computing and
classical (``regular'') computing research. That pipeline is summarized here because
it determines how papers entered the review queue in the first place, prior to the
individual manual review documented in Section~\ref{sec:results-detail}.

\subsubsection{Domain structure}
Both sides of the comparison organize their papers into six mirrored pipeline-stage
\emph{pillars}, so that each quantum research stage has a same-shape classical
counterpart (full topic and arXiv-category mappings are maintained centrally in a
domain-definitions module):

\begin{table}[htbp]
\centering
\caption{Mirrored pillar structure, quantum vs.\ classical computing}
\label{tab:pillars}
\begin{tabular}{@{}p{4.7cm} p{6.8cm}@{}}
\toprule
Quantum pillar & Classical (CS) mirror \\
\midrule
Hardware (superconducting qubits, transmons, fabrication) & Classical Hardware (CMOS/fabrication) \\
Control \& Operations & Classical Control \& Operations \\
Error Correction & Coding Theory \\
Algorithms & Classical Algorithms \\
Compilation \& Circuit Optimization & Compiler \& Circuit Optimization \\
Information Theory & Classical Information Theory \\
\bottomrule
\end{tabular}
\end{table}

Both sides of the comparison have all six pillars populated. On the quantum side, the
Hardware pillar is the one exception to the common construction procedure below: it
reuses an existing, separately-built dataset (Section~\ref{sec:quantum-list-gen}) built
before the six-pillar comparison architecture existed, rather than being queried by
the common procedure. Every other pillar on both sides --- including CS's own Hardware
pillar, which has no equivalent reused dataset and is queried the same way as every
other CS pillar --- is built by the common per-pillar arXiv query procedure
(Section~\ref{sec:pillar-corpus-construction}). This asymmetry is a known
inconsistency in the current pipeline, not a deliberate design choice; a planned
extension to resolve it is described in
Section~\ref{sec:planned-extensions}. The classical-computing pilot sample
(Table~\ref{tab:results-cs}) spans several of these pillars --- cryptography and
hardware architecture, compiler and circuit optimization, machine learning, and
algorithms --- while the quantum pilot sample (Table~\ref{tab:results-quantum}) is
drawn entirely from the Hardware pillar.

\subsubsection{Quantum list generation}
\label{sec:quantum-list-gen}

The superconducting-qubit corpus was built as a two-source pull, combining arXiv
preprints with journal publications not cross-listed on arXiv, then deduplicated into a
single list.

\paragraph{arXiv source.} Papers were retrieved from the arXiv API using a
title/abstract keyword match on \texttt{"superconducting qubit"} or
\texttt{"superconducting qubits"}, restricted to three arXiv categories
(\texttt{quant-ph}, \texttt{cond-mat.supr-con}, \texttt{cond-mat.mes-hall}), and
restricted to a submission-date window of 2021-01-01 through the present (a
\textasciitilde5-year window, chosen to keep the corpus current while remaining
tractable for a manual/LLM-assisted review pipeline). Results were paginated (100
records per request) and fetched with a politeness delay between requests per arXiv's
API guidelines.

\paragraph{Journal source (Crossref).} Because not every published paper in this area is
cross-listed on arXiv, a second pass queried the Crossref API with a bibliographic
search for \texttt{"superconducting qubit"}, restricted to the same date window and to
seven target journals: \emph{PRX Quantum}, \emph{npj Quantum Information},
\emph{Nature Physics}, \emph{Quantum}, \emph{Physical Review Letters}, \emph{Physical
Review A}, and \emph{Physical Review B}. Crossref's bibliographic search is a loose
relevance match rather than an exact-phrase filter, so the raw pull (1{,}880 hits, at
the most recent corpus rebuild) included a substantial fraction of off-topic
condensed-matter and general-physics results; a secondary title-relevance filter
(requiring the word ``qubit'' together with a superconducting-qubit-type term in the
title) removed 1{,}610 of these before merging.

\paragraph{Deduplication and merge.} The two sources were merged by normalized,
fuzzy-matched title comparison (a similarity threshold applied via sequence matching):
171 Crossref records matched an existing arXiv entry and were merged into that row
(adding venue/DOI metadata rather than creating a duplicate row); 99 Crossref records
had no arXiv match and were kept as journal-only rows. The resulting corpus totals
\textbf{1{,}485 papers} (1{,}386 arXiv-native $+$ 99 journal-only) as of the most recent
rebuild of this corpus; the sample reviewed in Section~\ref{sec:results-quantum} was
drawn from an earlier snapshot of this same pipeline (1{,}478 papers, 1{,}379
arXiv-native $+$ 99 journal-only), and the small difference between the two counts
reflects newly published papers accumulating between the two pulls rather than any
change in query or methodology.

\subsubsection{Pillar-based corpus construction (all pillars beyond Hardware)}
\label{sec:pillar-corpus-construction}

The five quantum pillars beyond Hardware, and all six CS pillars, are built by a
second, common procedure rather than the Crossref-augmented pipeline described above:
for every keyword defined for a pillar (topic and keyword lists are maintained
centrally in a domain-definitions module, one list per pillar per side), the arXiv API
is queried with a title/abstract exact-phrase match on that keyword, restricted to the
pillar's configured arXiv categories and the same 2021-present date window used
throughout this study, and the results for all of a pillar's keywords are merged and
deduplicated by arXiv identifier into one list per pillar; the six per-pillar lists on
a side are then merged and deduplicated again into that side's combined corpus.

\paragraph{Truncation defect and fix.} As originally implemented, each keyword query
retrieved only the first 100 results (arXiv's default single-page response), with no
pagination beyond that. This was found, during this study, to silently truncate any
keyword whose true match count exceeds 100 --- discovered by separately querying
arXiv's own reported \texttt{totalResults} count for every keyword on both sides,
which is not otherwise inspected by the pipeline. Several keywords were found to have
true totals in the thousands (e.g.\ ``quantum machine learning'': 1{,}668 true matches
against 100 previously retrieved; ``optimization algorithm'': 2{,}446 against 100), so
the original single-page implementation was retaining as little as 9--15\% of a
pillar's true candidate population in the worst-affected pillars (quantum
\emph{Algorithms} and \emph{Information Theory}; CS \emph{Algorithms}). The query
procedure was corrected to paginate every keyword query to its full reported total
rather than stopping at the first page, using the same page-size and politeness-delay
conventions already established elsewhere in this pipeline. This fix has been applied
to the CS-side pull (Section~\ref{sec:corpus-size}); the quantum-side pillars beyond
Hardware have not yet been re-pulled with the corrected procedure as of this writing.

\subsubsection{Keyword-list rework}
\label{sec:keyword-rework}

Correcting the truncation defect above exposed a second, independent problem that the
100-result cap had been masking: several CS keywords are generic, multi-word
descriptive phrases (e.g.\ \texttt{"machine learning algorithm"},
\texttt{"search algorithm complexity"}) that do not match how the literature they are
meant to capture actually names its own work. \texttt{"machine learning algorithm"}
returns 2{,}117 true matches, a small fraction of the true scale of the ML literature
on arXiv (well over 100{,}000 papers in \texttt{cs.LG} alone since 2021) --- not
because the pagination fix above failed, but because almost no machine-learning paper
literally contains that three-word phrase in its title or abstract; papers instead name
their specific method (``a transformer architecture for...'', ``a reinforcement
learning approach to...''). This is a phrasing problem, not a volume problem, and no
amount of additional pagination corrects it. By contrast, keyword lists built from
canonical, specific named techniques --- as most of the quantum-side list already is
(``Shor's algorithm'', ``QAOA'') and as CS's \emph{Coding Theory} pillar already is
(``LDPC code'', ``turbo code'', ``Reed-Solomon code'') --- do not exhibit this problem,
because papers in those subfields consistently name the specific technique they use.

The keyword lists on both sides were reworked on this basis: for pillars found to rely
on generic descriptive phrasing, specific named methods/techniques were added
alongside the existing keywords. No existing keyword was removed on either side ---
even a narrow, low-yield keyword still returns genuine matches, so removing it can only
lose valid coverage for no offsetting benefit. Table~\ref{tab:keyword-rework} lists
every change; all additions are purely additive over the lists in place when the
diagnostic pagination check above was run.

\begin{table}[htbp]
\centering
\caption{Keyword-list rework (additions only; no keyword was removed on either side)}
\label{tab:keyword-rework}
\begin{tabular}{@{}l l p{7.3cm}@{}}
\toprule
Side & Pillar & Keywords added \\
\midrule
Quantum & Information Theory & quantum Shannon theory, entanglement entropy \\
CS & Control \& Operations & model predictive control, adaptive control system \\
CS & Coding Theory & polar code, convolutional code \\
CS & Algorithms & convolutional neural network, transformer architecture, reinforcement
learning, graph neural network, generative adversarial network, dynamic programming,
approximation algorithm, combinatorial optimization \\
CS & Information Theory & Shannon capacity, differential privacy, zero-knowledge proof \\
\bottomrule
\end{tabular}
\end{table}

All other pillars on both sides (quantum Hardware, Control, Error Correction,
Algorithms, and Compilation; CS Hardware and Compiler \& Circuit Optimization) were
left unchanged: their existing keywords are already specific, canonical technique
names, and their true-total diagnostic did not show the generic-phrasing pattern
described above. As of this writing, the reworked keyword lists have been written into
the domain-definitions module but neither side's corpus has yet been re-pulled against
them; Section~\ref{sec:corpus-size} therefore still reports corpus sizes from before
this rework.

\subsubsection{Candidate list generation}
\label{sec:candidate-list-gen}

Once a domain's paper list exists (from either the quantum or the CS-side procedure
above), every paper in it is passed through a common data-availability screening step,
producing the \emph{candidate list} that manual review subsequently draws from. This
step is deliberately high-recall: per an explicit project requirement, the detector must
never silently classify a paper as having no data if any plausible signal is present.
False positives are an accepted cost of this design; false negatives are not.

\paragraph{Extraction.} For each paper, the PDF is downloaded from arXiv, its full
plain text is extracted (\texttt{pdftotext}), and its embedded hyperlink annotations
are extracted separately (\texttt{pdftohtml}, run with image extraction explicitly
disabled --- an early version of this step, run without that flag, silently dumped
every embedded figure as a separate image file, producing tens of thousands of stray
files before the fix). The PDF itself is deleted immediately afterward; only the
extracted text and link list are retained, since keeping the full-corpus PDFs was not
disk-feasible at this scale. Before pattern matching, the extracted text is
\emph{flattened}: hyphenated words and URLs split across a PDF line break are rejoined,
since an early version of the detector was found to miss valid links and DOIs broken
across a line wrap (e.g.\ \texttt{https:} at the end of one line, \texttt{//doi.org/...}
at the start of the next).

\paragraph{Signal detection.} Three independent signal sources are checked against the
flattened text: (i) a curated list of \textasciitilde30 data/code-availability
keyword phrases (e.g.\ \emph{data availab-}, \emph{code availab-}, \emph{available
upon}, \emph{deposited at}, \emph{openly available}, \emph{github repository}); (ii) a
list of \textasciitilde20 known data/code repository domains (GitHub, Zenodo, Figshare,
OSF, Dryad, institutional research-data repositories, etc.) checked against both PDF
hyperlink annotations and plain-text URLs; and (iii) generic DOI and URL regular
expressions scanned independently of nearby keywords, to catch a repository link that
appears without an explicit availability phrase nearby. A paper triggering any of these
three signals is labeled \emph{Candidate}; a paper triggering none is labeled \emph{No
candidate}. This is a binary triage, not a final classification --- it deliberately
does not attempt to distinguish a genuine self-citation (e.g.\ an author linking their
own prior repository) from a citation to an unrelated third-party tool, since that
distinction is left to the manual review step.

\paragraph{Corpus size and classification status at time of writing.}
\label{sec:corpus-size}

Table~\ref{tab:corpus-stats} reports the current pillar-corpus size on each side (all
six pillars combined) and how much of that corpus has completed the two-axis
classification described in Section~\ref{sec:extraction-classification}, as distinct
from the smaller, now-outdated candidate-count snapshot this table reported in earlier
drafts of this report.

\begin{table}[htbp]
\centering
\caption{Pillar-corpus size and classification status by domain, all six pillars}
\label{tab:corpus-stats}
\begin{tabular}{@{}l r r r@{}}
\toprule
Domain & Corpus size & Classified (Axis A/B) & Pending \\
\midrule
Superconducting qubits (quantum) & 3{,}478 & 3{,}472 (99.8\%) & 6 \\
Classical computing (CS, all six pillars) & 6{,}308 & 0 (0\%) & 6{,}308 \\
\bottomrule
\end{tabular}

\vspace{2mm}
\raggedright\footnotesize
The quantum figure reflects the pillar corpus as originally pulled, before either
corpus-construction fix in Sections~\ref{sec:pillar-corpus-construction}
and~\ref{sec:keyword-rework}; a re-pull incorporating both fixes is planned but had not
been executed as of this writing, and is expected to substantially increase this
corpus's size (a raw, pre-deduplication estimate places the corrected quantum corpus
at approximately 13{,}800--14{,}000 papers). Of the 6 quantum papers not yet
classified, all 6 failed PDF text extraction at the full-text-scan stage
(Section~\ref{sec:candidate-list-gen}) rather than being unclassifiable for any
framework-related reason; re-extraction has not yet been attempted. The CS figure
already reflects the pagination fix (Section~\ref{sec:pillar-corpus-construction}) but
not yet the keyword-list rework (Section~\ref{sec:keyword-rework}); classification of
this corpus had not yet begun as of this writing. Per-pillar corpus-size breakdowns for
both domains, and the true-total diagnostic table that motivated the two fixes above,
will be added to this report once the corrected, full-breadth corpora are pulled.

\end{table}

\subsection{Selection of papers for the detailed review in this report}
Two samples of fourteen papers each are individually reviewed in
Section~\ref{sec:results-detail}: a classical-computing sample selected from the CS-side
review queue described above, and a quantum-computing sample selected from the
superconducting-qubit review queue, chosen to span a comparable range of outcomes ---
directly-linked data, named-but-unlinked data, request-conditional access, and
unresolvable or not-yet-live links --- so that the two samples are a fair basis for
qualitative comparison despite their small size. Each paper, in both samples, was
processed independently as follows.

\subsection{Extraction, link verification, and classification}
\label{sec:extraction-classification}

Once a paper is selected from a domain's review queue (Section~\ref{sec:candidate-list-gen}),
it passes through three further stages before it appears in Table~\ref{tab:results-cs}
or Table~\ref{tab:results-quantum}. These stages are implemented as two scripts plus one
non-scripted classification step, described in turn below.

\paragraph{Stage 3: extraction and live link verification (scripted).} A dedicated
script reuses the plain text already cached by the candidate-generation step
(Section~\ref{sec:candidate-list-gen}) wherever available, avoiding a second download
of the same PDF; a paper not previously scanned is downloaded and extracted fresh using
the same procedure. Every link previously extracted for that paper is then
independently re-checked over HTTP (a HEAD request, falling back to a small ranged GET
for servers that reject HEAD), and classified as resolving, returning an error, or
timing out. This is the automated counterpart of the manual link-checking used earlier
in this project's pilot review --- the same kind of check that identified a claimed
GitHub repository returning a 404 and confirmed a code repository's stated pre-release
embargo. The result of this stage is one pending record per paper, containing the full
extracted text, the live link-verification results, and the classification fields left
unset.

This stage can be run in either of two scopes. The narrower scope processes only the
papers Stage 2 already flagged as Candidates; this is cheaper but inherits Stage 2's
blind spot (Section~\ref{sec:candidate-list-gen}) --- a paper that uses external data
without ever using an availability keyword or including a link is a Candidate-scope
false negative for Axis A. The broader scope processes every paper in the domain's
corpus regardless of its Stage 2 label, closing that gap at the cost of extracting and
verifying a larger set. For the quantum-side corpus, the broader scope was used, so
Stage 3's output covers all 3{,}478 papers across all six pillars, not only the subset
Stage 2 flagged.

\paragraph{Stage 4 batching (scripted preprocessing, not classification itself).} At
full-corpus scale, a paper's extracted text is too large to read in bulk without either
truncating it or exceeding practical per-request limits (a single paper's text ranges
from roughly 25{,}000 to over 500{,}000 characters, and the full quantum Stage 3 output
totals several hundred megabytes). Two scripted steps address this without changing what
is actually classified. First, the pending records for a batch of papers are condensed
into a digest: papers short enough to read whole are left as-is; longer papers are
reduced to their opening excerpt (title, abstract, and introduction, where a paper's
external data sources are typically first named), their closing excerpt (references and,
commonly, a formal data-availability statement), and a window of text around every
occurrence of a data/code-availability keyword found elsewhere in the paper, so a
relevant sentence in the middle of a long paper is not missed by only keeping the ends.
Classification is performed from this digest plus the unmodified link-verification
results, on the same basis it would be performed from the full text --- the digest
changes how much text is read, not what is asked of it or how it is judged. Second,
once a batch is classified, a separate script writes the classification fields onto a
copy of the paper's original, complete record (full extracted text and link-verification
results, unmodified), so the digest is never what persists in the domain's cache; it is
a reading aid discarded after use.

\paragraph{Stage 4: classification (not a scripted API call).} The two-axis, six-case
framework (Section~\ref{sec:framework}) is defined in exactly one place --- a single
framework document --- and every classification, whether performed by a person or by an
automated agent, is required to read that same document rather than work from a
separately maintained copy, so that the framework's definitions cannot drift out of
sync between where they are documented and how they are actually applied. Concretely,
classification in this project is performed by a Claude Code agent: for each pending
record produced by Stage 3, the agent reads the framework document together with that
paper's extracted text and link-verification results, and assigns an Axis A case, an
Axis B case, a short evidence quote for each, a confidence level, and a flag indicating
whether the assignment should be treated as provisional pending human confirmation. This
step is deliberately not implemented as a call to a hosted language-model API; the
extraction script that produces Stage 3's output is written so that it stops at
extraction and verification only, leaving classification to be performed and reviewed
directly rather than automated end-to-end without a human or agent in the loop at the
point of judgment.

\paragraph{Stage 5: review CSV construction (scripted).} A second script joins the
classification output back against the domain's paper-list metadata (title, authors,
publication date) and writes one CSV per domain, with human-readable case labels,
the evidence quotes, the link-verification outcomes, and a blank verification column.

\paragraph{Parallelization attempt and revision.} Batches were initially handed to
multiple independent classification agents running concurrently, each reading and
writing only its own batch's files (no two agents ever touch the same file, avoiding
write conflicts between them). In practice, all agents launched this way were
terminated by the same account-wide usage limit before any of them completed a batch,
so this parallel approach produced no usable output on the run attempted for this
project. Batches were subsequently processed sequentially instead, which trades
wall-clock time for the property that a failure partway through affects only the
batch in progress rather than all of them at once; work already merged from completed
batches is retained regardless of what happens to a later one.

\paragraph{Human verification.} No classification produced by Stage 4 is treated as
final on production. Every row is written with its verification column blank; a human
reviewer subsequently confirms or overrides each case assignment, following the same
convention already in use elsewhere in this project's review queues (verification
status recorded separately from the label itself, so the two are never conflated).
This mirrors the project's broader design principle, applied consistently from the
candidate-generation stage onward: an ambiguous or low-confidence signal is flagged for
review rather than silently resolved one way or the other.

\subsection{Domain (content) classification}
\label{sec:domain-classification}

Independently of the two-axis availability judgment above, every paper is also
assigned a primary subject pillar based on its actual content --- a separate
classification task from the pillar under which it was originally retrieved. Two
domain signals already exist earlier in this pipeline, and neither constitutes genuine
content classification: the pillar recorded at corpus-construction time
(Section~\ref{sec:pillar-corpus-construction}) reflects only which keyword search
happened to retrieve a paper, and the full-text scan (Section~\ref{sec:candidate-list-gen})
additionally checks whether any of a pillar's keyword strings appear anywhere in a
paper's full text, which is still substring presence rather than judgment --- a paper
can accumulate several such tags purely from background citations without being
substantively about any of them. A paper's true subject can genuinely differ from
both: a hardware paper that mentions ``optimization algorithm'' only in passing should
not be tagged \emph{Algorithms} on that basis alone.

Domain classification is performed with the same reading depth and per-paper judgment
as the availability classification (Section~\ref{sec:framework}), not a reduced or
heuristic pass: an agent reads a paper's extracted text (or digest, at full-corpus
scale, per Section~\ref{sec:extraction-classification}) and independently judges which
one of the six mirrored pillars (Table~\ref{tab:pillars}) the paper is actually about,
rather than inferring this from keyword presence.

The judgment space for this task is not limited to the six pillars: a paper can also
be classified as out of scope for quantum computing entirely. This matters
specifically on the quantum side, where several pillar keywords (e.g.\ ``quantum
entanglement,'' ``quantum communication,'' ``quantum measurement theory'') are used as
heavily in non-computing quantum physics --- quantum sensing, quantum foundations,
quantum optics, quantum networking --- as in quantum computing itself, and the
corpus-construction procedure (Section~\ref{sec:pillar-corpus-construction}) is
deliberately kept broad (maximizing recall, per the same design principle already
applied to the full-text availability scan in Section~\ref{sec:candidate-list-gen})
rather than pre-filtered for this distinction at pull time. Rather than narrowing the
query and losing volume, this distinction is deferred to domain classification, reusing
the out-of-scope category set already established earlier in this project for exactly
this problem: \emph{Atomic Physics}, \emph{Quantum Sensing}, \emph{Quantum
Magnonics}, \emph{Quantum Optomechanics}, \emph{Photonic Quantum Computing},
\emph{Trapped-Ion Computing}, \emph{Neutral-Atom Computing}, \emph{Semiconductor Spin
Qubits}, and \emph{High Energy Physics}, plus a general \emph{Not Quantum Computing}
fallback for a paper that does not fit any of the above. A paper assigned one of these
labels is retained in the corpus (so that corpus size and candidate-generation yield
remain auditable) but excluded from every subsequent quantum-computing-specific
statistic, including the pillar-level results this report is ultimately built around.
No equivalent out-of-scope category is currently defined on the CS side; the same
question --- whether CS's broad pull similarly needs an out-of-scope carve-out --- is
tracked as an open item in Section~\ref{sec:planned-extensions}.

For papers not yet classified for
availability, this judgment is obtained in the same pass as the Axis A/Axis B
classification --- the same read of the paper's text yields both judgments, avoiding a
second full read-through of the corpus. For the 3{,}472 quantum papers already
classified for availability before this domain-classification methodology was adopted
(Table~\ref{tab:corpus-stats}), a separate backfill pass adds only the primary-domain
judgment, at the same rigor as if it had been obtained in the combined pass; their
existing Axis A/Axis B classification, evidence, and confidence fields are left
unmodified. This backfill is implemented as its own reusable script, applying the same
deterministic-extraction-plus-agent-classification pattern already established for
availability classification, so that domain classification can be independently
re-run in the future (e.g.\ following a revision to the pillar definitions in
Table~\ref{tab:pillars}) without requiring the availability judgment to be redone. As
of this writing, this domain-classification methodology has been designed and the
supporting keyword-list corrections (Section~\ref{sec:keyword-rework}) have been made,
but neither the backfill pass over the 3{,}472 already-classified quantum papers nor
the combined-pass classification of the CS corpus has yet been executed.

\subsection{Scope note}
The samples reported in Section~\ref{sec:results-detail} were classified through direct
manual review, predating the Stage 3--5 pipeline described in
Section~\ref{sec:extraction-classification}; that pipeline is intended for scaling
classification to the full corpus in each domain (Table~\ref{tab:corpus-stats}) going
forward, following the same procedure and the same framework document applied by hand
here. As of this writing, Stage 3 has been run against the entire quantum-side corpus
(all 3{,}478 papers across all six pillars, not only the subset Stage 2 flagged as
Candidates), and Stage 4 classification of that output is complete for 3{,}472 of
those 3{,}478 papers (Table~\ref{tab:corpus-stats}) via the batched procedure above;
this is not yet reflected in Table~\ref{tab:results-quantum}, which remains the
earlier hand-reviewed pilot sample, and the classified full-corpus results are not yet
included in this report (see the note at the end of the abstract). This is, in either case, a single-pass
review that has not yet been
cross-validated by a second reviewer, and one paper in the working sample (a
Tsetlin-Machine-based load-monitoring paper) was reviewed for domain classification only
and is omitted from the results table pending a full data-availability pass.

\section{Classification Framework}
\label{sec:framework}

\subsection{Axes}

\textbf{Axis A --- External Data.} Whether the paper draws on data that already existed
independently of the paper --- a pre-collected dataset, real-world records, an
established benchmark, existing designs or artifacts --- and performs its analysis on
that.

\textbf{Axis B --- Self-Generated Data.} Whether the paper produces its own data by
running an experiment, simulation, or generative procedure --- random instance
generation, physics-based simulation, hardware simulation, training-run outputs, etc.
--- and performs its analysis on that.

A paper is scored independently on each axis; it may use both, either, or neither. A
paper is classified as purely theoretical only when \emph{both} axes independently
return Case 1.

\subsection{Cases (applied identically to both axes)}

\begin{table}[htbp]
\centering
\caption{Case definitions}
\label{tab:cases}
\begin{tabular}{@{}>{\bfseries}p{2.6cm} p{10.5cm}@{}}
\toprule
Case & Definition \\
\midrule
Case 1 --- Not present &
This axis is not applicable to the paper. For Axis A: no external data was used at
all. For Axis B: no self-generated data was produced at all. \\
\addlinespace
Case 2 --- Present, directly accessible &
The data is named and accompanied by a working access path (URL, DOI, or repository)
that lets an independent party retrieve the identical material. Ideally version- or
commit-pinned. \\
\addlinespace
Case 3 --- Present, not handed over &
The data is named, described, or cited, but no access path is given. An independent
party must locate it themselves using only the name or description provided. \\
\addlinespace
Case 4 --- Present, underspecified &
The data is acknowledged to exist but described too vaguely or incompletely to
independently identify or reconstruct it --- e.g., missing dataset name/version (Axis
A), or missing generation parameters/methodology (Axis B). \\
\addlinespace
Case 5 --- Claimed accessible, unverifiable &
The paper explicitly states the data is available and gives a specific access path,
but that path does not resolve. Distinct from Case 3 in that access was affirmatively
promised and failed to hold. \\
\addlinespace
Case 6 --- Conditional on request &
The paper states the data will be provided if the authors are contacted directly,
without any immediate, self-service access path. Distinct from Case 3 (no offer at
all) and Case 5 (no present-availability claim). \\
\bottomrule
\end{tabular}
\end{table}

\section{Results}
\label{sec:results-detail}

Table~\ref{tab:results-cs} reports the Axis A and Axis B classification obtained for
each classical-computing paper reviewed in this pass; Table~\ref{tab:results-quantum}
(Section~\ref{sec:results-quantum}) reports the same for the matched quantum-computing
sample.

\subsection{Classical computing sample}

\begin{longtable}{@{}p{4.3cm} p{1.3cm} c c p{5.3cm}@{}}
\caption{Per-paper data availability classification --- classical computing sample} \label{tab:results-cs} \\
\toprule
Paper & arXiv ID & Axis A & Axis B & Note \\
\midrule
\endfirsthead
\toprule
Paper & arXiv ID & Axis A & Axis B & Note \\
\midrule
\endhead
\bottomrule
\endfoot

Matrix multiplication exponent (AlphaEvolve) & 2608.16884 & 1 & 1 &
Pure mathematical optimization over a fixed analytical object; no dataset of any
kind. \\
\addlinespace
MIS/Max-K-SAT convergence & 2608.18910 & 1 & 2 &
Random graphs/formulas generated by the authors' own code; generator openly linked on
GitHub. \\
\addlinespace
HQC on ARM Cortex-M4 & 2608.18145 & 1 & 1 &
No dataset; evaluation is cycle-count benchmarking on hardware. Standard KAT test
vectors are used only for correctness checking, not analysis. \\
\addlinespace
Communication-reduction DMPC via LSTM & 2608.17592 & 1 & 3 &
No external dataset; training data is 2{,}000+ simulated multi-robot scenarios with
parameters specified in the text, but no code/data repository was located. \\
\addlinespace
Convex losses for SVM/SVR/shallow NN & 2608.14288 & 3 & 1 &
Seven UCI Machine Learning Repository datasets named and tabulated; no direct link
given in the paper text. \\
\addlinespace
SynAct (LLM agent for logic synthesis) & 2608.12751 & 3 & 1 &
14 open-source RTL designs from OpenCores named as the source; no direct link or
version given. \\
\addlinespace
ARMOR (RTL simulation acceleration) & 2608.08462 & 3 & 1 &
Designs generated via TensorLib/Chipyard/RocketChip/BOOM/Gemmini, named but not
linked; the paper's own code repository exists but states it is embargoed until after
MICRO 2026. \\
\addlinespace
Fraudulent banking operations recognition & 2608.07471 & 2 & 1 &
Credit Card Fraud Detection dataset (Kaggle, ULB) with a working direct URL given in
an explicit Data Availability Statement. \\
\addlinespace
Synthetic LiDAR data generation & 2608.07106 & 3 & 2 &
Base CAD models drawn from the existing ModelNet dataset (named, not directly linked
in text); the physics-simulated LiDAR dataset the authors derive from it is
separately released with a working Zenodo DOI. \\
\addlinespace
MiCoPro (mixed-precision HW/SW co-design) & 2608.06916 & 3 & 1 &
Six standard ML benchmarks (CIFAR-10/100, MNIST, ImageNet, Speech Commands V2,
TinyStories) named and cited by reference number, not linked; the authors' own
framework code is separately linked on GitHub. \\
\addlinespace
PDMarks (physical-design IP watermarking) & 2608.05055 & 2 & 1 &
Eight open-source RTL benchmarks drawn from OpenROAD-related repositories, linked
with pinned commit hashes. \\
\addlinespace
Generalized Quadratic Gradient & 2608.01552 & 3 & 5 &
Several real datasets named only in figure captions, with no citation or link in the
body text; the paper's own experiment code is explicitly claimed to be ``openly
available'' at a GitHub URL that was verified to return a 404. \\
\addlinespace
Signum-based distributed optimization/ML & 2608.01220 & 2 & 6 &
MNIST (external) is linked directly via Kaggle in a Data Availability Statement; the
authors' own raw simulation output (from a self-generated synthetic regression
example) is stated to be available only ``upon request.'' \\
\addlinespace
Domain-adaptive Deep JSCC & 2607.28907 & 3 & 1 &
SVHN, MNIST, and PACS named and cited by reference number, not linked in text; the
authors' own code repository is real and functional, and auto-downloads MNIST only ---
SVHN and PACS still require the reader to independently source and stage the raw
files. \\

\end{longtable}

\subsection{Quantum computing sample}
\label{sec:results-quantum}

Table~\ref{tab:results-quantum} reports the same classification for a matched sample of
fourteen papers drawn from the superconducting-qubit review queue
(Section~\ref{sec:quantum-list-gen}).

\begin{longtable}{@{}p{4.3cm} p{1.3cm} c c p{5.3cm}@{}}
\caption{Per-paper data availability classification --- quantum computing sample} \label{tab:results-quantum} \\
\toprule
Paper & arXiv ID & Axis A & Axis B & Note \\
\midrule
\endfirsthead
\toprule
Paper & arXiv ID & Axis A & Axis B & Note \\
\midrule
\endhead
\bottomrule
\endfoot

Flip-chip integrated SC qubits & 2608.07306 & 1 & 2 &
No external dataset; own fabrication/measurement data and processing code both
directly linked (4TU data repository DOI and GitHub, cited as numbered references). \\
\addlinespace
Component-level inverse design of transmon qubits & 2607.20795 & 3 & 2 &
Training data derived from the external SQuADDS database, named but not directly
linked in this paper's own text; the authors' own processed splits, scalers, and
trained model weights are directly available in a linked GitHub repository. \\
\addlinespace
Radiopurity assays for SC qubits at SNOLAB & 2607.16151 & 1 & 2 &
No external dataset used as an analysis input; the paper's own newly-measured
material-radioactivity assay data is deposited in and directly accessible from the
community database radiopurity.org (verified live). \\
\addlinespace
Real-time charge-jump detection (CNN) & 2607.14293 & 1 & 3 &
No external dataset; own training data (Ramsey tomography scans from NEXUS
measurements) is described in the text but no access path of any kind is given. \\
\addlinespace
Kibble-Zurek excitations on a 5,564-qubit annealer & 2607.13842 & 1 & 2 &
No external dataset; own quantum-annealer and matrix-product-state simulation data
directly deposited on Zenodo with a working DOI. \\
\addlinespace
Plaquette (FTQC hardware-aware design platform) & 2607.08767 & 1 & 6 &
No external dataset identified in the available text; the paper's own data is stated
to be available only ``from the authors upon reasonable request'' -- the platform
itself is separately noted to be a commercial product with its own restricted access
terms. \\
\addlinespace
Pauli-propagation noise-canceling observables & 2606.20441 & 1 & 3 &
No explicit data-availability statement of any kind located; two GitHub links present
are official Qiskit-organization repositories rather than a confirmed own-data
release, so self-generated data (if any) is treated as described-but-unconfirmed. \\
\addlinespace
Non-perturbative CPMG scaling (EmuPlat) & 2603.29525 & 1 & 2 &
No external dataset; own simulation data and scripts stated to be available now at a
GitHub repository (the base repository resolves; the specific cited sub-path returns
404). A more permanent archival copy is separately stated to be deposited ``upon
acceptance,'' i.e.\ not yet live at time of review. \\
\addlinespace
Charge-parity detection via randomized benchmarking & 2604.02809 & 1 & 2 &
No external dataset; own source data files directly linked via a working Figshare DOI,
with additional (unspecified) data available only upon request -- a partial release
rather than a complete one. \\
\addlinespace
Lattice-renormalized tunneling models for qubit materials & 2512.18156 & 1 & 5 &
No external dataset; own simulation data explicitly claimed available, but the cited
Dryad DOI is a literal placeholder (\texttt{10.5061/dryad.XXXXXXXXX}) that does not
resolve to any record. \\
\addlinespace
Magnon squeezing in the quantum regime & 2602.19671 & 1 & 2 &
No external dataset; own source data for the figures directly deposited on Zenodo
(verified live), with additional information available only upon request -- partial
but the core figure data is genuinely open. \\
\addlinespace
Resist-free shadow deposition for JJ fabrication & 2604.09796 & 1 & 2 &
No external dataset; own fabrication and measurement data openly available at a
working Zenodo DOI. \\
\addlinespace
LLM-assisted superconducting qubit experiments (HAL) & 2603.08801 & 1 & 2 &
No external dataset; the paper's own control-software framework (HAL) and its
supporting tools (Grapher, QuICK) are all self-citations to the same lab's own prior
open-source releases, all independently verifiable on GitHub -- though specific
experimental-result data from the demonstration itself is not separately addressed. \\
\addlinespace
Scalable open-source QEC decoding system & 2603.16203 & 1 & 3 &
No external dataset; the paper is framed as ``open-source'' and a plausibly-own
repository (EosCore) was located, but no explicit data/code availability statement
confirming this was found in the reviewed text. \\

\end{longtable}

\section{Summary of Observations}

\subsection{Classical computing sample}

\begin{itemize}[leftmargin=1.5em]
    \item \textbf{Case 3 (``named, not handed over'') is the modal outcome for Axis
    A} in this sample: six of the papers that use external data cite it by name only,
    leaving retrieval entirely to the reader.
    \item \textbf{Only two papers achieve Case 2 on Axis A with the strongest form of
    reproducibility} --- a direct, working link. Of these, only one (PDMarks) pins the
    reference to an exact commit hash, guaranteeing a byte-identical copy regardless of
    future changes to the upstream repository.
    \item \textbf{Two papers make an explicit availability claim that does not hold up
    under verification} (Case 5) or resolves only to a conditional promise rather than
    an immediate path (Case 6). Both were only detected as such because the stated
    access paths were independently checked rather than taken at face value ---
    underscoring that a stated ``available at ...'' claim is not equivalent to verified
    availability.
    \item \textbf{Self-generated data (Axis B) is comparatively rare} in this sample
    and, when present, is unevenly documented: one paper links its generator code
    directly (Case 2), one paper describes its generation procedure but links nothing
    (Case 3), and two papers leave their own generated output only conditionally or
    unverifiably accessible (Cases 5 and 6).
    \item \textbf{No paper in this sample returned Case 1 on both axes except the
    purely mathematical optimization paper}, consistent with the expectation that
    Case-1/Case-1 should be reserved for papers with no empirical component at all.
\end{itemize}

\subsection{Quantum computing sample}

\begin{itemize}[leftmargin=1.5em]
    \item \textbf{Axis A is overwhelmingly Case 1 in this sample} --- thirteen of the
    fourteen papers use no external dataset at all; the sole exception (the
    transmon-design paper) names its external source (SQuADDS) without a direct link in
    the reviewed text. This is the sharpest structural difference from the classical
    sample, where external-data use (Cases 2--6 on Axis A) was the norm rather than the
    exception.
    \item \textbf{Axis B (self-generated data) is where nearly all the variation in this
    sample lives}, since almost every paper reports its own newly fabricated,
    measured, or simulated data. Eight of fourteen papers reach Case 2 (directly
    accessible), consistent with a real, if inconsistently applied, community norm of
    depositing experimental data to Zenodo or an institutional repository under a
    formal ``Data availability'' heading.
    \item \textbf{Partial releases (data linked, with additional data gated behind a
    request) appear three times} (charge-parity detection, magnon squeezing, and,
    differently, the CPMG-scaling paper's pending archival deposit) --- distinct from
    the classical sample, where a request-only outcome (Case 6) was typically total
    rather than partial.
    \item \textbf{One paper's stated access path is a literal unassigned placeholder}
    (\texttt{10.5061/dryad.XXXXXXXXX}) rather than a broken but once-valid link, a
    distinct failure mode from the classical sample's 404'd GitHub URL: here the paper
    was evidently submitted before the archival deposit was even created, not merely
    before its retrieval was verified.
    \item \textbf{Self-citation to a lab's own prior open-source tooling} (as opposed to
    releasing new data specific to the paper at hand) appears in this sample (the
    LLM-assisted-experiments paper, citing its own group's HAL/Grapher/QuICK releases)
    and was also observed repeatedly in the broader superconducting-qubit review queue
    beyond this fourteen-paper sample --- a citation style that reads, on first pass, like
    a third-party tool reference and requires checking the citing lab's identity against
    the repository owner to correctly attribute.
\end{itemize}

\subsection{Cross-domain comparison}

Read together, the two samples suggest that the two computing communities structure
their data-availability problem differently rather than simply exhibiting different
overall care. The classical sample's difficulty is overwhelmingly an Axis A problem:
papers routinely build on named, citable, but unlinked external benchmarks and
datasets. The quantum sample's difficulty is instead concentrated in Axis B: since these
papers almost always generate their own new experimental or simulated data, the
practice being tested is not ``did you cite your source'' but ``did you actually deposit
what you produced,'' and the sample shows a genuine, if partial and unevenly applied,
norm of doing so (Case 2 on Axis B in 8/14 papers). A same-shape comparison across the
full review queues, once both are fully verified, is needed before this observation can
be treated as more than suggestive --- both samples reported here are small,
non-random, hand-selected subsets chosen for outcome diversity rather than
representativeness (see Limitations).

\section{Planned Extensions}
\label{sec:planned-extensions}

This section tracks methodology changes that have been decided on but not yet
executed, identified while reviewing this report's description of the pipeline against
what the pipeline should actually do. It is expected to grow as further review surfaces
further such directions; each entry stays here until the corresponding pipeline change
has been made and the rest of this report updated to reflect it.

\subsection{Unify Hardware-pillar corpus construction with the common procedure}
\label{sec:planned-hardware-unification}

As noted in Section~\ref{sec:candidate-generation}, the quantum-side Hardware pillar is
currently built by a separate, older two-source procedure (arXiv plus Crossref
cross-referencing against seven target journals) rather than the common per-pillar
arXiv query procedure used by every other pillar on both sides --- an inconsistency
left over from before the six-pillar comparison architecture existed, not a deliberate
design choice. The two-source method is, if anything, more thorough than the common
procedure (it additionally catches journal-only publications with no arXiv preprint),
so the fix is to extend it outward rather than to drop it: the Crossref
cross-referencing step will be added to the common per-pillar procedure
(Section~\ref{sec:pillar-corpus-construction}) for every pillar on both sides, using
each pillar's own keyword list and an appropriately chosen set of target journals,
rather than rebuilding the Hardware pillar under the weaker arXiv-only method it was
originally exempted from.

This extension will be applied even to pillars, and specifically to the quantum
Hardware pillar, that already have completed classification work
(Table~\ref{tab:corpus-stats}). Any additional papers the Crossref cross-referencing
step finds beyond what is already in a pillar's corpus will be deduplicated against the
existing corpus by arXiv identifier (and, for journal-only records with no arXiv
identifier, by the same fuzzy title-matching procedure already used in
Section~\ref{sec:quantum-list-gen}) and added to it; no paper already in a corpus, and
no classification already completed against it, will be removed or discarded as a
result of this extension. This follows the same additive-only principle already
applied to the keyword-list rework (Section~\ref{sec:keyword-rework}): a corpus-widening
change should only ever add coverage, never retract it.

\subsection{Decide whether CS needs an out-of-scope category analogous to quantum's}
\label{sec:planned-cs-out-of-scope}

Section~\ref{sec:domain-classification} resolves the quantum side's
computing-vs-physics scope question by classifying out-of-scope papers into an
existing set of non-computing labels (\emph{Quantum Sensing}, \emph{Atomic Physics},
etc.) at the domain-classification stage, rather than restricting the corpus-pull
query and losing recall. Whether the CS side has an analogous problem --- a keyword
matching a paper that is not really about the intended pillar's subject at all, as
opposed to a paper that is on-topic but assigned to the wrong one of the six pillars
--- has not yet been assessed. A candidate example: ``channel capacity'' and
``information theoretic limit'' (CS info\_theory) are both used extensively in pure
information-theory mathematics and in communications-engineering work with no
connection to computing. This needs the same kind of review the quantum side already
received before deciding whether CS's domain classifier needs its own out-of-scope
label set, and if so, what it should contain.

\section{Limitations}

This pass reflects a single reviewer's manual reading of each paper and has not been
cross-checked by a second coder; category boundaries (particularly between Case 3 and
Case 4) required judgment calls that a second reviewer might resolve differently.
Link verification was performed once, at the time of review, and some URLs may change
status afterward. The Tsetlin-Machine load-monitoring paper in the working CS sample was
not carried through the full data-availability procedure and is therefore omitted from
Table~\ref{tab:results-cs} rather than assigned a placeholder classification. Both the
classical and quantum samples reported here (fourteen papers each) are small,
non-random subsets of their respective review queues (380 CS candidates and 235
quantum candidates), hand-selected for outcome diversity rather than drawn at random,
and should not be treated as representative of those queues, of the six pillars on
either side as a whole, or of the quantum-vs-classical comparison this pipeline is
ultimately intended to support. The quantum sample is drawn entirely from the Hardware
pillar; the cross-domain comparison in Section~\ref{sec:results-quantum} should
accordingly be read as Hardware-pillar-only and may not generalize to the other five
quantum pillars.

Beyond the pilot-sample limitations above, three pieces of work described in this
report are designed but not yet executed, and every corpus size and classification
count given here should be read with that in mind. First, the quantum-side pillar
corpus has not yet been re-pulled with the pagination fix and keyword-list rework
described in Sections~\ref{sec:pillar-corpus-construction}
and~\ref{sec:keyword-rework}; the 3{,}478-paper figure and the 3{,}472 papers already
classified against it (Table~\ref{tab:corpus-stats}) both predate that correction, and
the corrected corpus is expected to be several times larger (an approximate
raw estimate places it at 13{,}800--14{,}000 papers). Second, the CS-side corpus
(6{,}308 papers) reflects the pagination fix but not the keyword-list rework, and has
not yet been through any stage of classification. Third, the domain-classification
methodology (Section~\ref{sec:domain-classification}) has been designed but neither
its backfill pass over the already-classified quantum papers nor its combined-pass
application to new classification work has been run. None of the numbers in
Table~\ref{tab:corpus-stats} or Section~\ref{sec:results-detail} should be treated as
final until these three steps are complete and reflected in an updated version of this
report, which is also where the full-corpus results tables and the corpus-construction
diagnostic tables (Table~\ref{tab:keyword-rework} and the underlying true-total
analysis) will be added.

\end{document}
